\section{生物学示例族 IV：双稳态与序列机制}

许多神经系统与行为系统都呈现双稳态或序列机制结构。
双稳态之所以重要，是因为同一个系统可能支持两个相对稳定、并由阈值分开的结果。
序列机制之所以重要，是因为顺序本身会成为系统架构的一部分。

本项目已索引的那篇关于自组织吸引子的论文，在这里尤其有价值。
它的中等强度贡献，不仅在于吸引子的存在，
更在于它说明：重复暴露于序列，会重塑耦合结构，并稳定后续的序列行为。

\begin{example}[把学习序列看作被学会的转变图]
设一位读者反复执行同一序列：
\begin{enumerate}[nosep]
  \item 坐在同一张桌前，
  \item 打开同一本笔记，
  \item 回顾上一次未解决的那一行，
  \item 以一道例题开局。
\end{enumerate}
随着时间推移，这个序列本身就会变得更容易被进入并被走完。
序列动力学语言的数学价值在于：目标不再只是某一个状态，
它也可以是一条穿过状态空间的\emph{路径}。
\end{example}